Die zuverlässige Bereitstellung und der kontinuierliche Betrieb von Softwareanwendungen 
bilden heute einen fundamentalen Baustein moderner IT-Infrastrukturen. Während die 
Entwicklung von Backend- und Frontend-Komponenten nach wie vor im Mittelpunkt steht, 
rückt parallel dazu die Frage in den Vordergrund, wie Anwendungen verlässlich, 
wiederholbar und mit vertretbarem Wartungsaufwand betrieben werden können. Diesem 
Themenkomplex widmet sich das vorliegende Kapitel, indem es die Containerisierung 
der entwickelten Anwendung sowie deren Bereitstellung auf einer virtuellen Maschine 
beleuchtet.

Innerhalb des Gesamtprojekts werden fachliche und funktionale Themen, darunter die 
Umsetzung von Backend und Frontend sowie die Implementierung von Authentifizierungs- 
und Autorisierungsmechanismen, von anderen Teammitgliedern bearbeitet. Dieses Kapitel 
konzentriert sich dagegen gezielt auf die technische Infrastruktur, die den Betrieb 
der Anwendung trägt. Im Zentrum steht dabei die Fragestellung, wie durch Container-
Technologien eine einheitliche Ausführungsumgebung etabliert werden kann, die 
plattformunabhängig funktioniert.

In den vergangenen Jahren hat sich Containerisierung als weitverbreiteter Standard 
für die Bereitstellung von Software etabliert. Technologien wie Docker erlauben es, 
Anwendungen mitsamt sämtlicher Abhängigkeiten in isolierten Containern zu bündeln. 
Dadurch werden klassische Herausforderungen der Softwareverteilung, etwa inkompatible 
Bibliotheksversionen oder komplexe manuelle Konfigurationsschritte, spürbar 
reduziert \cite{pahl2015containerization}. Für das hier beschriebene Projekt ergibt 
sich daraus ein wesentlicher Nutzen: Die Anwendung läuft sowohl in der Entwicklungs- 
als auch in der Produktivumgebung unter vergleichbaren Bedingungen.

Das Kapitel verfolgt das Ziel, die verwendeten Container-Technologien und den 
konkreten Deployment-Ablauf transparent zu dokumentieren. Es wird dargelegt, wie 
Docker zur Containerisierung genutzt wird, wie Docker Compose das koordinierte 
Zusammenspiel mehrerer Dienste ermöglicht und wie die Bereitstellung auf einer 
virtuellen Maschine erfolgt. Ein besonderes Augenmerk liegt dabei auf der 
Nachvollziehbarkeit der technischen Entscheidungen, um künftige Wartungsarbeiten 
und Erweiterungen zu vereinfachen.

Zur inhaltlichen Eingrenzung ist festzuhalten, dass erweiterte Infrastrukturthemen 
wie Continuous Integration, Continuous Deployment oder Cloud-native Plattformen 
nicht behandelt werden. Kubernetes wird zwar im weiteren Verlauf konzeptionell 
eingeordnet, kommt jedoch nicht zum produktiven Einsatz. Der Schwerpunkt liegt 
bewusst auf einer überschaubaren, gut handhabbaren Lösung, die den Projektanforderungen 
gerecht wird und einen praxisnahen Betrieb auf einer virtuellen Maschine 
ermöglicht \cite{dockerdocs}.